\chapter{Evaluation Case Studies}
\label{app:evaluation-cases}

The evaluation chapter summarizes the batch results. This appendix gives more
detail on representative cases that shaped the final design. The cases are not
included to claim that every streaming site behaves the same way. They show the
kind of failures that appeared during testing and the engineering response used
to make the workflow more stable.

\section{Case Study A: Landing Page With Repeated Event Cards}

Some websites expose the first useful evidence on a discovery page rather than
on a player page. The page lists matches, leagues, or channels, and each event
may link to one or more downstream hosting pages. The first prototype often
stopped after collecting a few visible links because it treated the landing page
as a normal link list.

\begin{table}[!htbp]
\centering
\begin{tabularx}{\textwidth}{Q{3.8cm}Y}
\toprule
\rowcolor{accent!8}
Evaluation point & Observation \\
\midrule
Initial symptom & The agent returned only the first visible match cards and
missed links hidden under date tabs or lower page sections. \\
Main risk & Host-candidate recall looked acceptable on small pages but failed
when the same layout repeated by date, league, or event block. \\
Change introduced & The landing agent was required to inspect visible sections,
open date or category tabs, scroll when needed, and keep a ledger of already
checked blocks. \\
Acceptance signal & The run produced a ranked list of hosting candidates with
the associated match context and no duplicate navigation noise. \\
\bottomrule
\end{tabularx}
\caption{Landing-page regression case: repeated event cards.}
\label{tab:case-landing-repeated-cards}
\end{table}

This case changed the way landing pages were evaluated. A run could not be
counted as working simply because it found one event. It had to show that the
visible discovery surface was inspected deliberately. That is why host-candidate
recall became a separate KPI in Chapter~6.

\section{Case Study B: Hosting Page With Popup Interference}

Hosting pages were usually harder than landing pages because evidence appears
only after interaction. The agent has to click play, switch servers, close
popups, restore focus, and check whether the player state changed. In early
runs, the agent often clicked a play button, landed on an advertisement tab, and
continued reasoning from the wrong page.

\begin{table}[!htbp]
\centering
\begin{tabularx}{\textwidth}{Q{3.8cm}Y}
\toprule
\rowcolor{accent!8}
Evaluation point & Observation \\
\midrule
Initial symptom & The browser opened a popup or new tab after the first click,
and the agent interpreted the advertisement page as the current investigation
target. \\
Main risk & The final trace looked active because tools were still executing,
but the evidence path had left the original player page. \\
Change introduced & After every interaction, the browser layer checks for new
contexts, closes unrelated popups, restores focus to the original page, and
retries the click with a safer strategy. \\
Acceptance signal & The screenshot and media-state check both refer to the
original hosting page or to a controlled embedded fallback, not to an
advertising tab. \\
\bottomrule
\end{tabularx}
\caption{Hosting-page regression case: popup and tab focus recovery.}
\label{tab:case-hosting-popup}
\end{table}

This case explains why screenshot evidence is not optional. A network result
alone does not prove that the agent stayed on the correct page. The screenshot
shows whether the player page, popup, overlay, or blocker was visible at the
moment the claim was made.

\section{Case Study C: Embedded Player With Delayed Manifest}

Embedded pages sometimes look simple because they contain only one player.
However, the stream manifest may not be requested until after a play action,
after a short delay, or after an iframe redirect. Early harvesting often
returned no media evidence because the tool ran before the player produced the
network requests.

\begin{table}[!htbp]
\centering
\begin{tabularx}{\textwidth}{Q{3.8cm}Y}
\toprule
\rowcolor{accent!8}
Evaluation point & Observation \\
\midrule
Initial symptom & The embedded agent called the harvest tool immediately and
returned an empty stream list. \\
Main risk & The site was incorrectly marked as unsupported even though the
manifest appeared after player activation. \\
Change introduced & The embedded agent now inspects player readiness, clicks
play when needed, waits for network activity, captures a screenshot, and only
then harvests media evidence. \\
Acceptance signal & The run returns a manifest or media request after an
observed player-state change, or it records a clear failure reason if activation
is blocked. \\
\bottomrule
\end{tabularx}
\caption{Embedded-player regression case: delayed manifest generation.}
\label{tab:case-embedded-delayed-manifest}
\end{table}

The lesson from this case was that tool order matters as much as tool
availability. The same harvest tool can be ineffective or useful depending on
when it is called. This is why the final prompt requires observe--act--verify
steps before extraction.

\section{Case Study D: Tokenized Stream URL}

Some players generate stream URLs that include short-lived tokens. The URL can
work during the run and become invalid later. This behavior created two
requirements: preserve the URL exactly as observed, and preserve enough context
to understand where it came from.

\begin{table}[!htbp]
\centering
\begin{tabularx}{\textwidth}{Q{3.8cm}Y}
\toprule
\rowcolor{accent!8}
Evaluation point & Observation \\
\midrule
Initial symptom & The run found a media URL, but a later manual replay failed
because the token had expired. \\
Main risk & A reviewer could mistake token expiry for a false extraction result. \\
Change introduced & The evidence bundle stores the URL, timestamp, page URL,
frame source, screenshot, provider lookup, and failure context. \\
Acceptance signal & A tokenized stream is accepted only with trace context:
where it was observed, when it was captured, and which page state produced it. \\
\bottomrule
\end{tabularx}
\caption{Stream-evidence case: tokenized URL preservation.}
\label{tab:case-tokenized-url}
\end{table}

This case also influenced provider enrichment. A tokenized URL should be
enriched quickly while its host and path are still visible in the run context.
Even if the token later expires, the provider and network evidence remain useful
for review.

\section{Case Study E: Wrong Page Classification}

Classification failures were usually not caused by the model misunderstanding a
single word. They were caused by pages that mixed several roles: a landing page
with an embedded player, a hosting page with a match list, or an error page with
old player markup still present in the DOM.

\begin{table}[!htbp]
\centering
\begin{tabularx}{\textwidth}{Q{3.8cm}Y}
\toprule
\rowcolor{accent!8}
Evaluation point & Observation \\
\midrule
Initial symptom & The classifier routed a mixed page to the wrong specialist,
which caused the next agent to use unsuitable tools. \\
Main risk & A landing agent may collect too many irrelevant links from a player
page, while a hosting agent may click controls that do not exist on a listing
page. \\
Change introduced & Classification now considers visible page role, player
signals, iframe context, and action needed, rather than URL shape alone. \\
Acceptance signal & The chosen route gives the next specialist a task that can
produce useful evidence or a defensible failure reason. \\
\bottomrule
\end{tabularx}
\caption{Routing case: mixed-role page classification.}
\label{tab:case-classification-mixed-role}
\end{table}

This case justified the handoff object in the architecture chapter. The
orchestrator does not only pass a URL. It passes the page role, route source,
evidence checklist, candidate context, and navigation policy.

\section{Case Study F: Provider Enrichment Without Abuse Contact}

Provider enrichment was useful even when it did not return an abuse contact.
Some hosts resolve to a provider or network owner, but the lookup path does not
expose a direct abuse email. Earlier reports treated that as a failure. The
final evaluation separates provider coverage from abuse-contact discovery.

\begin{table}[!htbp]
\centering
\begin{tabularx}{\textwidth}{Q{3.8cm}Y}
\toprule
\rowcolor{accent!8}
Evaluation point & Observation \\
\midrule
Initial symptom & A stream URL resolved to an IP address and provider, but no
abuse email was returned. \\
Main risk & The evidence enrichment looked like a total failure even though the
provider and network context were useful. \\
Change introduced & The dashboard and evaluation split checked URLs, resolved
IPs, provider matches, and abuse contacts into separate metrics. \\
Acceptance signal & Provider coverage is counted when provider or network
context is found; abuse-contact discovery remains a separate downstream field. \\
\bottomrule
\end{tabularx}
\caption{Provider-enrichment case: useful metadata without abuse contact.}
\label{tab:case-provider-no-abuse}
\end{table}

This distinction makes the evaluation more honest. The project can show that
provider enrichment works for most extracted URLs while still admitting that
abuse-contact discovery needs additional sources and manual review.

\section{Case Study Summary}

\begin{longtable}{>{\raggedright\arraybackslash\bfseries}p{3.4cm}
                  >{\raggedright\arraybackslash}p{4.7cm}
                  >{\raggedright\arraybackslash}p{5.5cm}}
\caption{Summary of evaluation case studies and design impact.}
\label{tab:case-study-summary}\\
\toprule
\rowcolor{accent!8}
Case family & Failure pattern & Design response \\
\midrule
\endfirsthead
\toprule
\rowcolor{accent!8}
Case family & Failure pattern & Design response \\
\midrule
\endhead
\bottomrule
\endfoot
Landing discovery & Visible cards were collected, but hidden date tabs or lower
sections were missed. & Landing-specific inspect, checked-section ledger, and
host-candidate recall metric. \\
Hosting interaction & Popup or ad tab stole focus after a click. & Popup
closing, focus restoration, screenshot checkpoint, and safer click strategies. \\
Embedded extraction & Harvesting ran before the player emitted media requests.
& Activation guard, media-state check, and delayed harvest. \\
Tokenized streams & URL became invalid after the run. & Timestamped stream
record, screenshot context, provider enrichment, and trace preservation. \\
Mixed-role pages & Classifier chose the wrong specialist. & Route based on page
role and required action, not only URL pattern. \\
Provider enrichment & Provider found, but abuse contact missing. & Separate
metrics for checked URLs, resolved IPs, provider matches, and abuse contacts. \\
\end{longtable}

These cases explain why the final evaluation is regression-oriented. Each fix
improved one behavior, but every fix also had to be checked against pages that
already worked. The case studies therefore connect implementation decisions to
the evaluation metrics used in the main report.

\section{Manual Review Checklist for Case Studies}

For each representative case, the manual review followed the same questions.
This made the case studies comparable and avoided accepting a run because it
looked successful at first glance.

\begin{itemize}
  \item Was the initial page role visible in the screenshot or DOM signals?
  \item Did the selected specialist match the page role and evidence objective?
  \item Did the trace show the browser action that changed the page state?
  \item Did screenshots support the final claim or the failure reason?
  \item Was stream evidence captured after player activation rather than before
  it?
  \item If provider enrichment failed, was the missing field explicit?
  \item Did the fix improve this case without breaking older successful cases?
\end{itemize}

The checklist is intentionally simple. Its purpose is to keep evaluation
grounded in evidence that a reviewer can inspect: page role, action, state
change, stream result, provider context, and regression impact.
